\documentclass[11pt]{article}
\usepackage{icmlko}
\begin{document}

\title{접전에서는 아직 동전이다 --- 판의 난이도로 가르지 않으면 우리 점수는 아무 말도 하지 않는다}
\author{월드모델 실험실}
\date{노트 527 · 2026년 8월 4일}
\maketitle

\begin{abstract}
``쌍 정확도 $67.4\%$'' 는 기준선이 없으면 뜻이 없다는 지적을 받고 두 가지를
쟀다. ① \textbf{실무자가 쓸 법한 규칙}과 견주기 ② \textbf{판의 난이도}로
가르기.

\begin{center}
\begin{tabular}{lrrr}
\hline
기준 & 전체 & \textbf{접전} & 명백 \\
\hline
\textbf{우리 모델} & $\mathbf{67.4\%}$ & $\mathbf{56.1\%}$ & $\mathbf{81.7\%}$ \\
도메인별 최고 축(학습 선택) & $64.4\%$ & $55.3\%$ & $73.5\%$ \\
공통 규칙 하나(학습 선택) & $58.7\%$ & $53.0\%$ & $65.6\%$ \\
\hline
\end{tabular}
\end{center}

접전 $=$ 실제 성적 백분위 차가 $20$ 점 미만인 쌍($284{,}053$ 쌍),
명백 $=$ $50$ 점 이상($198{,}370$ 쌍).

\textbf{이기기는 이긴다}(단일 축 대비 $+3.0$ \%p, 공통 규칙 대비 $+8.7$ \%p).
그러나 \textbf{접전에서는 $56.1\%$ --- 거의 동전이고 단일 축보다 $+0.8$ \%p
나을 뿐이다.} 우리가 버는 곳은 \textbf{명백한 판}($+8.1$ \%p)이다.
\end{abstract}

\begin{figure}[t]
\centering
\includegraphics[width=\textwidth]{figs/f1.pdf}
\caption{\textbf{가} --- 난이도별. \textbf{나} --- 도메인별 우리 대 최고 축.
\textbf{다} --- 읽기.}
\end{figure}

\section{먼저, 내 첫 기준선이 틀렸다}

처음에는 \emph{유보 표본 위에서} 축 $70$ 개(정방향 $\cdot$ 역방향)를 전부 돌려
1등을 골랐다. 그 값은 $68.0\%$ 로 \textbf{우리 모델보다 높았다.} 그런데 그건
사람이 낼 수 있는 수가 아니라 \textbf{채점표를 보고 고른 신탁}이다.

규칙을 \textbf{학습 표본에서} 고르고 유보에서 채점하는 것으로 고치니
$64.4\%$ 가 됐다. \textbf{$3.6$ \%p 가 사후 선택의 값이었다}(노트 133 의
``첫 수를 그대로 받지 않는다''가 이번에는 \emph{기준선} 쪽에서 걸렸다).

\begin{quote}
\textbf{도서가 표본이다} --- 학습에서 $62.6\%$ 로 1등이던 규칙(\texttt{grp})이
유보에서 \textbf{$36.9\%$}, \emph{동전보다 나쁘다}. ``지난번에 통했던 규칙''이
어떻게 배신하는지가 한 줄에 있다. 우리 모델은 같은 도메인에서 $63.6\%$ 다.
\end{quote}

\section{접전과 명백은 다른 문제다}

\begin{center}
\begin{tabular}{lrrr}
\hline
도메인 & 접전 & 명백 & 격차 \\
\hline
모바일 & $58.2\%$ & $90.0\%$ & $31.8$ \\
게임 & $58.0\%$ & $89.5\%$ & $31.5$ \\
애니 & $57.8\%$ & $88.5\%$ & $30.7$ \\
아이돌 & $52.9\%$ & $82.5\%$ & $29.6$ \\
\hline
\end{tabular}
\end{center}

\begin{quote}
\textbf{실무자는 명백한 판을 이미 맞힌다.} 그러니 지금 우리 값어치는
``명백한 판을 \emph{자동으로 그리고 규모로} 처리하는 것''이고,
\textbf{접전은 아직 아무도 못 판다.}
\end{quote}

\textbf{왜 접전이 안 되는지도 같은 자료가 말한다.} 지금 모델의 입력은
\emph{기획서 속성}뿐이다. 접전을 가르는 것은 시점 $\cdot$ 경쟁작 $\cdot$ 운영
$\cdot$ 그때 그 사람들이고, 그건 \textbf{상태 $\cdot$ 개인 $\cdot$ 개입 층}에
있다. 로드맵이 그 순서인 이유가 여기서 처음으로 \emph{측정}으로 나왔다.

\section{축 하나가 지배하는 도메인에서는 진다}

\begin{center}
\begin{tabular}{lrrl}
\hline
도메인 & 우리 & 최고 축 & 그 축 \\
\hline
시장팝업 & $58.0\%$ & $\mathbf{71.1\%}$ & \texttt{goods\_scale} \\
아이돌 & $66.7\%$ & $\mathbf{78.2\%}$ & \texttt{venue\_prominence} \\
모바일 & $72.5\%$ & $\mathbf{83.2\%}$ & \texttt{entry\_friction} \\
팝업 & $60.2\%$ & $\mathbf{67.1\%}$ & \texttt{target\_breadth} \\
\hline
\end{tabular}
\end{center}

넷에서 진다. 반대쪽 끝은 도서 $+26.7$ \%p $\cdot$ 펀딩 $+9.2$ $\cdot$ 웹툰
$+8.1$ 이다.

\begin{quote}
\textbf{읽기(가설)} --- 한 도메인에 \emph{지배적인 축}이 있으면 공유 모델이
그것을 \textbf{희석한다.} \texttt{entry\_friction} 은 도메인마다 부호가
뒤집혀(노트 441) 집 밖 예보에서 가리고 있는 축인데, 모바일 안에서는 단독으로
$83.2\%$ 다. 공유 모델이 다른 도메인의 반대 부호에 끌려 모바일에서 이 축을
덜 쓰고 있을 수 있다.
\end{quote}

\textbf{다음 관문(사전 등록)} --- \texttt{DOMDROP} 의 거울, 곧 \textbf{도메인별
축 강조}를 시험한다. 노트 525 가 가림 계열을 닫았지만 \emph{강조}는 아직 안
쟀다. 예측: 지는 넷에서 오르고 나머지는 안 변한다. 틀릴 수 있는 방식 ---
이웃이 그만큼 잃으면 판은 안 오른다(노트 525 에서 섞기가 그랬다).

\section*{아직 안 잰 것 --- 사람 기준선}

\textbf{우리는 사람을 한 번도 안 쟀다.} 단일 축은 \emph{대리}일 뿐이다.
설계는 이렇게 적어 둔다.

\begin{itemize}
\item 유보 표본에서 \textbf{층화 추출한 쌍 $100$ 개}(접전 $50$ $\cdot$ 명백 $50$)를
실무자에게 보여 주고 고르게 한다. 모델과 \emph{같은 정보}만 준다(기획서 속성).
\item 응답자 $\ge 5$ 명, 도메인당 최소 $10$ 쌍. 순서 무작위화.
\item \textbf{판정}: 실무자 중앙값이 접전에서 $56.1\%$ 를 넘으면 우리는 접전에서
사람보다 못한 것이고, 그 사실을 그대로 적는다.
\end{itemize}

\section*{남은 것}

\begin{itemize}
\item \textbf{도메인별 축 강조}(위 사전 등록) --- 지는 넷이 후보다.
\item \textbf{접전 층에서만 채점하는 자}를 관문에 넣을지 --- 지금 다섯 자는
모두 전체 쌍이라 접전 개선을 못 본다.
\item \textbf{사람 기준선} --- 사람이 필요하다.
\end{itemize}

\end{document}
